% !TEX root = ../neurips_2026.tex

\section{Background}
\label{sec:background}

We introduce notation and recall the concepts that our theoretical analysis builds on: autoregressive knowledge distillation, input saliency, and the Jacobian of the input--output map.

\paragraph{Autoregressive knowledge distillation.}
Let $f_T$ and $f_S$ denote the teacher and student language models.
Given an input sequence $x = (x_1, \ldots, x_L)$ with embeddings $e = (e_1, \ldots, e_L) \in \R^{L \times d}$, each model produces a next-token distribution $p_M(\cdot \mid x, t) \in \Delta^{|\mathcal{V}|-1}$ at every position $t$, where $\mathcal{V}$ is the vocabulary.
We partition positions into \emph{prompt} positions $\mathcal{P} = \{1, \ldots, L_p\}$ and \emph{response} positions $\cR = \{L_p{+}1, \ldots, L\}$.
Standard KD~\citep{hinton2015distilling} trains the student to minimize the forward KL divergence over response tokens:
\begin{equation}\label{eq:standard-kd}
  \cL_{\mathrm{KD}}
  = \frac{1}{n} \sum_{i=1}^{n} \underbrace{\sum_{t \in \cR_i} D_{\kl}\!\bigl( p_T(\cdot \mid x_i, t) \;\big\|\; p_S(\cdot \mid x_i, t) \bigr)}_{D(x_i) \;\triangleq\; \text{per-sample KL}},
\end{equation}
where distributions are temperature-scaled by $T > 1$ and the loss is multiplied by $T^2$ following standard practice~\citep{hinton2015distilling}.
We write $D(x) = \sum_{t \in \cR} D_t(x)$ for the per-sample KL summed over response positions, with $D_t(x) = \sum_{v \in \mathcal{V}} p_T^v(x,t) \log \frac{p_T^v(x,t)}{p_S^v(x,t)}$.

\paragraph{Input saliency.}
The \emph{input saliency} at position $j$ measures how much the model's response generation depends on the $j$-th input token~\citep{wu-etal-2023-ad}:
\begin{equation}\label{eq:saliency-def}
  s_j(x) \;=\; \left\| \frac{\partial \log P(\text{response} \mid x)}{\partial e_j} \right\|_2,
  \qquad
  \log P(\text{response} \mid x) = \sum_{t \in \cR} \log p(x_t \mid x_{<t}).
\end{equation}
The saliency vector $\mathbf{s}(x) = (s_1, \ldots, s_{L_p}, 0, \ldots, 0) \in \R^L$ is nonzero only at prompt positions: it summarizes \emph{which input tokens the model relies on} to generate its response.
We write $\mathbf{s}_T(x)$ and $\mathbf{s}_S(x)$ for teacher and student saliency vectors.
For sample reweighting (Section~\ref{sec:reweighting}), we convert saliency to a distribution via temperature-scaled softmax over prompt positions: $\hat{s}_j = \exp(s_j / \tau_s) / \sum_{k \in \mathcal{P}} \exp(s_k / \tau_s)$.

\paragraph{Input--output Jacobian.}
The Jacobian of model $M$'s output with respect to input embeddings is $J_M(x) \in \R^{V \times (L \cdot d)}$, where entry $(v, (j{-}1)d + k)$ gives $\partial [\log p_M^v] / \partial e_{j,k}$.
This matrix decomposes into $L$ blocks $J_M = [J_M^{(1)}, \ldots, J_M^{(L)}]$ where $J_M^{(j)} \in \R^{V \times d}$ captures the sensitivity of all output logits to the $j$-th input embedding.
A key structural observation connects the Jacobian to saliency: the per-position saliency is the Frobenius norm of the corresponding Jacobian block,
\begin{equation}\label{eq:saliency-jacobian}
  s_j(x) = \|J_M^{(j)}(x)\|_F,
\end{equation}
since $\|\partial \log P / \partial e_j\|_2 = \|J_M^{(j)}\|_F$ when summing over the response log-probability (we formalize this in Lemma~\ref{lem:saliency-jacobian}).
Thus saliency compresses each $V \times d$ Jacobian block into a single scalar, preserving the per-position magnitude while discarding directional information within each block.

\paragraph{Sobolev norms and function approximation quality.}
In classical approximation theory, the $L^2$ norm $\|f_T - f_S\|_{L^2} = (\int |f_T(x) - f_S(x)|^2 \, dx)^{1/2}$ measures pointwise agreement, while the Sobolev $W^{1,2}$ norm additionally penalizes derivative mismatch:
$\|f_T - f_S\|_{W^{1,2}}^2 = \|f_T - f_S\|_{L^2}^2 + \|\nabla f_T - \nabla f_S\|_{L^2}^2$.
By the Sobolev embedding theorem, $W^{1,2}$ convergence implies uniform convergence in sufficiently regular spaces---an approximation that is close in $W^{1,2}$ must also be close in every neighborhood of every training point.
\citet{czarnecki2017sobolev} brought this perspective to neural network training, showing that matching derivatives (``Sobolev training'') accelerates learning and improves generalization.
Our work applies this principle to knowledge distillation: standard KD is $L^2$ approximation (output matching); adding Jacobian matching via noise perturbation upgrades it to $W^{1,2}$ (output + sensitivity matching).

\paragraph{Existing distillation methods.}
Table~\ref{tab:methods} summarizes the landscape of LLM distillation methods.
All operate at the zero-order level: they vary \emph{which} divergence to minimize between teacher and student outputs, but none constrains \emph{how} the student processes inputs.
MiniLLM~\citep{gu2024minillm} reverses the KL direction;
GKD~\citep{agarwal2024gkd} uses Jensen--Shannon divergence with on-policy sampling;
DistiLLM~\citep{ko2024distillm} employs skew-KL for stability;
SeqKD~\citep{kim2016sequence} trains on teacher-generated sequences;
DA-KD~\citep{dakd2025} adaptively reweights samples by output-level difficulty.
The only methods that incorporate first-order information---AD-KD~\citep{wu-etal-2023-ad} and gradient KD~\citep{wang2022gkd}---are restricted to encoder-only BERT models for classification, requiring explicit gradient computation that is infeasible for decoder-only LLM generation.

\begin{table}[t]
\centering
\caption{Comparison of distillation methods by approximation order, Jacobian handling, and applicability. Only \sagd{} achieves implicit first-order matching for autoregressive LLM generation.}
\label{tab:methods}
\small
\begin{tabular}{@{}lcccc@{}}
\toprule
\textbf{Method} & \textbf{Order} & \textbf{Jacobian signal} & \textbf{Reweighting} & \textbf{Setting} \\
\midrule
Standard KD~\citep{hinton2015distilling} & 0th & --- & uniform & any \\
MiniLLM~\citep{gu2024minillm} & 0th & --- & uniform & LLM \\
SeqKD~\citep{kim2016sequence} & 0th & --- & uniform & LLM \\
GKD~\citep{agarwal2024gkd} & 0th & --- & uniform & LLM \\
DistiLLM~\citep{ko2024distillm} & 0th & --- & uniform & LLM \\
DA-KD~\citep{dakd2025} & 0th & --- & output KL & LLM \\
\midrule
AD-KD~\citep{wu-etal-2023-ad} & 1st & IG attribution (MSE) & uniform & BERT cls.\ \\
Gradient KD~\citep{wang2022gkd} & 1st & full gradient (MSE) & uniform & BERT cls.\ \\
Sobolev Training~\citep{czarnecki2017sobolev} & 1st & full Jacobian & --- & small nets \\
Jacobian Matching~\citep{srinivas2018knowledge} & 1st & full Jacobian & --- & CNN cls.\ \\
\midrule
\textbf{\sagd{} (ours)} & \textbf{0th+1st} & \textbf{implicit (noise KL)} & \textbf{saliency JSD} & \textbf{LLM gen.} \\
\bottomrule
\end{tabular}
\end{table}
